% Original book diagram: fig:Posterior_Contraction
% Book source: sections/Part_IV/bayesian_learning.tex
% Original diagram authorship retained. Visual design: Ziad Mosalam.
\documentclass[11pt,border=8pt]{standalone}
\usepackage{amsmath,amssymb,graphicx,tikz}
\usetikzlibrary{arrows.meta,automata,positioning,calc,shapes.geometric,fit,backgrounds,patterns}
\newcommand{\vect}[1]{\mathbf{#1}}
\newcommand{\mat}[1]{\mathbf{#1}}
\newcommand{\Normal}{\mathcal{N}}
\newcommand{\prob}{p}
\newcommand{\E}{\mathbb{E}}
\newcommand{\Var}{\mathrm{Var}}
\newcommand{\Cov}{\mathrm{Cov}}
\newcommand{\CovP}{\mathbf{P}}
\newcommand{\Prob}{\mathbb{P}}
\setlength{\textwidth}{6.5in}
% Shared visual vocabulary for the book's original diagrams.
% Figures opt in explicitly, so unrelated text and existing additions are unchanged.
\definecolor{BookBlue}{HTML}{003262}
\definecolor{BookGold}{HTML}{B17A00}
\definecolor{BookTeal}{HTML}{247C86}
\definecolor{BookInk}{HTML}{30363B}
\definecolor{BookLight}{HTML}{E9EFF5}
\tikzset{
  book figure/.style={font=\small\rmfamily, color=BookBlue, text=BookInk,
    line width=0.8pt, >=Latex},
  book axis/.style={->, draw=BookInk, line width=0.6pt},
  book link/.style={->, draw=BookBlue, line width=0.8pt},
  book observation/.style={draw=BookGold, fill=BookGold!8},
  book state/.style={draw=BookBlue, fill=BookLight},
  book note/.style={font=\footnotesize\rmfamily, text=BookInk}
}

\pdfinfoomitdate=1
\pdftrailerid{}
\pdfsuppressptexinfo=15
\begin{document}
\begin{tikzpicture}[book figure, scale=1.5]
\draw[book axis] (-0.2,0)--(6,0) node[right]{$\theta$};
\draw[book axis] (0,-0.1)--(0,2.2);

% Gaussian posteriors with standard deviation proportional to 1/sqrt(n); areas are equal
\draw[BookTeal, book curve, domain=0.2:5.8, samples=120] plot (\x,{0.40*exp(-0.5*((\x-3.25)/1.00)^2)});
\draw[BookGold, book curve, dashed, domain=1.2:4.8, samples=120] plot (\x,{0.89*exp(-0.5*((\x-3.1)/0.45)^2)});
\draw[BookBlue, book curve, dotted, domain=2.0:4.0, samples=120] plot (\x,{1.81*exp(-0.5*((\x-3.03)/0.22)^2)});
\draw[BookInk!60, dashed] (3,0) -- (3,2.05) node[above, BookInk] {\small $\theta_{\text{true}}$};
\node[text=BookTeal] at (4.75,0.45){$n=1$};
\node[text=BookGold] at (3.95,0.85){$n=5$};
\node[text=BookBlue] at (3.55,1.75){$n=20$};
\end{tikzpicture}
\end{document}
